\scanpage{25}
$\mathbf C$, on a un objet canonique \deleted{de} $\xi$ de $T_{\mathbf C}(\mathbf Z)$
[$T$ le torseur fondamental associé à $\mathbf C$).

\diagram{\begin{tikzcd}[ampersand replacement=\&]H^i(X,\mathbf Q(0)) \arrow[r,rightsquigarrow] \& H^i(X_{\mathrm{top}},\mathbf Z)\end{tikzcd}}

L'hom $H^i(X^n,\mathbf Q(0))\longrightarrow H^{\unsure{i}}(Y^{\unsure{m}},\mathbf Q(\unsure{0}))$ défini par $\deleted{\unsure{z^n}}\ z\in A^n(X\times Y)$ est celui défini \inserted{par voie trans., à l'aide de} \deleted{par la classe} $P_{\mathrm{top}}(z)\left[=\dfrac{1}{(2i\pi)^n}P_{\mathrm{alg}}(z)\right]$
\inserted{\[
\boxed{\xi(z)(x)=\operatorname{pr}_{2*}^{\unsure{\mathrm{an}|\mathrm{top}}}\bigl(\operatorname{pr}_1^*(x)P_{\mathrm{top}}(z)\bigr)}
\]}

Considérons l'\unsure{isom} can (pour $X,i$ donnés)
\diagram{\begin{tikzcd}[ampersand replacement=\&,column sep=huge,row sep=large]
 \xi_{\mathbf C}(H^{\unsure{i|p}}(X,\mathbf Q(\unsure{\infty|0})))=\xi(H^{\unsure{i|p}}(X,\mathbf Q(\unsure{\infty|0})))\otimes_{\mathbf Z}\mathbf C \arrow[r,"\sim"] \& \deleted{\unsure{\underline{\underline H}^*}}\underline{\underline H}^i(X,\Omega_X^*)=\xi_{\mathbf C,\infty}(H^{\unsure{i|p}}(X,\mathbf Q(\unsure{\infty|0}))) \\
 H^{\unsure{i|p}}(X_{\unsure{(h)}},\unsure{\mathbf C}) \arrow[u,equal] \arrow[ru,dash,"{\sim\;\deleted{\mathrm{can}}}","{\text{can par \unsure{Hodge-Gro.}}}"'] \&
\end{tikzcd}}
Cet isom \uline{\deleted{\unsure{est}} \unsure{est} \deleted{\unsure{donc}} fonctoriel} en $H^p(X,\mathbf Q(\unsure{\infty|0}))$.

\deleted{Pour le rendre canonique, il faut}
\inserted{En effet, si $z\in A^n(X\times Y)$, on a}
\[
\begin{aligned}
 x\rightsquigarrow\operatorname{pr}_{2*}^{\mathrm{top}}\bigl(\operatorname{pr}_1^*(x)P_{\mathrm{top}}(z)\bigr)
 &=\operatorname{pr}_{2*}^{\mathrm{top}}\left(\operatorname{pr}_1^{\unsure{a|*}}(x)\left(\frac1{2i\pi}\right)^nP_{\mathrm{alg}}(z)\right)\\
 &=\left(\frac1{2i\pi}\right)^{\unsure{-n|-m}}\operatorname{pr}_{2*}^{\mathrm{alg}}\left(\operatorname{pr}_1^*(x)\left(\frac1{2i\pi}\right)^nP_{\mathrm{alg}}(z)\right)
\end{aligned}
\]
\diagram{\begin{tikzcd}[ampersand replacement=\&,row sep=huge,column sep=huge]
 H^p(X,\mathbf C) \arrow[r,"{\xi_{\mathbf C}(z)\;\deleted{\unsure{\mathrm{alg}}}}","{\substack{\text{défini par}\\P_{\mathrm{top}}(z)\\=\left(\frac1{2i\pi}\right)^nP_{\mathrm{alg}}(z)}}"'] \arrow[d,dash,"\wr\ \mathrm{can}"] \& H^p(Y,\mathbf C) \arrow[d,dash,"\wr\ \mathrm{can}"] \\
 \underline{\underline H}^p(X,\Omega_X^\bullet) \arrow[r,"{\substack{\xi_{\mathbf C,\infty}^{\unsure{\mathrm{alg}}}(z)\\\text{défini par }P_{\mathrm{alg}}(z)\\=\left(\frac1{2i\pi}\right)^{\unsure{n|m}}\xi_{\mathbf C}(z)}}"'] \& H^p(Y,\Omega_X^\bullet)
\end{tikzcd}}
\diagram{\begin{tikzcd}[ampersand replacement=\&]x\arrow[r,rightsquigarrow] \& \operatorname{pr}_{2*}^{\mathrm{alg}}\bigl(\operatorname{pr}_1^*(x)\underbrace{P_{\mathrm{alg}}(z)}\bigr)\end{tikzcd}}

\figure{\begin{tikzpicture}\draw (0,0)--(.6,2.3);\draw (.08,.35)--(.65,2.3);\end{tikzpicture}}
\margin{Cette compatibilité peut déjà \unsure{s'énoncer} pour analytique $\longleftrightarrow$ \unsure{topologique} (\unsure{avec classes de corr. analytiques})}

\margin{Ouf !}
\mspage{24}
